\documentclass[12pt, titlepage]{article}
%\usepackage[norsk]{babel}
\usepackage[utf8]{inputenc}
\usepackage{minted}
\usepackage{pdflscape, diagbox}
\usepackage[colorlinks = true,
linkcolor = blue,
urlcolor = blue]{hyperref}


\begin{document}

\title{Mappeoppgave - del 1 \\ Tillegsinformasjon - Vurdering}
\author{IIRA2001 - Høst 2023}
\date{}
\maketitle

\section*{Vurdering}

\textit{Utkast til sensurveiledning av kodebiten til mappe 1 følger.  Veiledning til Refleksjonsnotat kommer senere}

\subsection*{Kodeferdigheter}

\subsubsection*{E}

Du har klart å lage pythonprogrammer som løser oppgavene vi har jobbet med i timen, slik den er beskrevet i oppaveteksten.
Fremgangsmåten tilfredstiller et minimum, men har et ferdighetsmessig forbedringspotensiale i feks å:
\begin{itemize}
	\item {bruke fornuftige deskriptive variabelnavn}
	\item {legge gjenbrukbare deler av programmet i egne funksjoner}
	\item {strukturere data i egnede strukturer som lister}
	\item {legge repeterende kode i løkker}
\end{itemize}
Du viser i liten grad vurderingsevne og selvstendighet i praktisk bruk av programmene.

\subsubsection*{D}
Du har klart å lage pythonprogrammer som løser oppgavene vi har jobbet med i timen.
Framgangsmåten din har mangler, men viser bruk av flere av konseptene vi har jobbet med, slik som:
\begin{itemize}
	\item {å bruke fornuftige deskriptive variabelnavn}
	\item {å legge gjenbrukbare deler av programmet i egne funksjoner}
	\item {å strukturere data i egnede strukturer som lister}
	\item {å legge repeterende kode i løkker}
\end{itemize}
Du viser nokså god vurderingsevne og selvstendighet ved å bruke noen av programmene til noe konkret og tolke eller diskutere resultatene.
Du har for eksempel kanskje:
\begin{itemize}
	\item {Brukt andre meningsfulle parametre enn de i oppgaveteksten fra timene og grunngitt valgene}
	\item {Gjort mindre forandringer i programmet slik at det løser et litt annerledes problem}
	\item {På andre måter vist vurderingsevne/selvstendighet i problemstilling eller diskusjon}
\end{itemize}

\subsubsection*{C}

Du har laget pythonprogrammer som omhandler problemstillinger nært knyttet til oppgavene fra timen.
Koden bruker programmeringskonseptene vi har jobbet med på en god måte, og du viser god vurderingsevne og selvstendighet.
Programmene dine skiller seg ut fra oppgavebeskrivelsene fra timen ved kanskje å:
\begin{itemize}
	\item {Løse en ganske annerledes problemstilling enn de gitt i oppgavene fra timen}
	\item {Ha utvidet funksjonalitet}
\end{itemize}

\textit{Å vise vurderingsevne og selvstendighet kan dere til eksempel gjøre ved å bruke programmet deres til å løse et konkret problem,
	eller å undersøke et konkret tilfelle dere selv finner på. Det vi ønsker å se er at dere kan bruke
	python som verktøy i en (mer eller mindre) faglig sammenheng. Sett også av litt plass til å diskutere resultatene; hva fant du ut?}


\subsubsection*{B}
Du har laget programmet som er ganske løsrevet fra slik som de er beskrevet i oppgavebeskrivelsene fra timen.
Du viser gode programmeringsferdigheter og selvstendighet ved å anvende programmering til noe konkret.
For en B ser vi spesielt etter at du kan anvende programmeringsverktøy i nye sammenhenger, altså konkrete problem vi ikke har gitt i timen.

\subsubsection*{A}
For en A ser vi etter meget høy grad av selvstendighet og evne til å bruke programmering i nye sammenhenger.
Du har kanskje brukt python på en utforskende måte til å gi faglig innsikt

\subsection*{Total vurdering}

Dere skal levere 3 programmer - det er ikke slik at alle programmene må vurderes til en 'A' for at det skal bli endelig vurdering.
Man kan legge mer arbeid i ett av programmene, og gjør de 2 andre litt enklere i utførelsen.



%\begin{landscape}
%	\begin{table}
%		\begin{center}
%			\small
%			\begin{tabular}{c| c| c| c| c| c| c| }
%				\hline
%				\backslashbox{Kode}{Refleksjonsnotat} & Viser innsikt i bla & Ser bla bla & Tjohei & ltrlsldkjfk & dkljsdf & lsdkjflkjj \\
%				\hline
%				Løser en oppgave greit                & A                   & B           & C      & D           & E       & F          \\
%				Bruker programmet til noe konkret     & A                   & B           & C      & D           & E       & F          \\
%				Bruker programmet                     & A                   & B           & C      & D           & E       & F          \\
%				Løser en oppgave greit                & A                   & B           & C      & D           & E       & F          \\
%				Løser en oppgave greit                & A                   & B           & C      & D           & E       & F          \\
%			\end{tabular}
%		\end{center}
%	\end{table}
%\end{landscape}

\section*{Refleksjonsnotat}

Ytterligere vurderingsveiledning/info kommer senere

\section*{Samarbeid/Plagiat}

Det er naturlig at dere samarbeider og hjelper hverandre når dere skal løse problemer og lære å kode.
Når en programmererer er det også veldig vanlig å ta bruk av eksempelkoder fra nettet, for eksempel fra dokumentasjon, tutorials eller stack-overflow.

Vi er derfor ikke like strenge på plagiatkontroll/kilder for python-koden.  
Dersom du og en medelev har samarbeidet tett, kan dere oppgi navnene til hverandre. 
Prøv likevel å skille oppgavene deres litt ved for eksempel å forskjellig input som gir vesentlig annerledes resultat, om mulig

Dersom dere har hentet litt komplisert kode fra nettet som dere ikke har full kontroll på og ikke har forandret særlig selv,
oppgi gjerne dette samt hvor dere har hentet koden fra.

Det viktigste er at diskusjon og refleksjon er et individuelt arbeid.



\section*{FAQ}

\subsubsection*{Hvordan skal vi levere i BlackBoard?}
\textit{Informasjon om innlevering er nylig lagt til beskrivelsen på innleveringen i blackboard}

\subsubsection*{ ...}


\end{document}
